% SPDX-License-Identifier: CC-BY-SA-4.0
% Copyright (c) 2026 Rafael Beraldo
% Parte 5 (fim): operação no grupo, com o mundo, repetidoras e protocolo
% Incluído por presentation.tex; compile a partir dele.

\begin{frame}{Comunicação entre o grupo ou a cordada}
  \textbf{Simplex}: comunicação direta, na mesma frequência, sem repetidora. Só funciona se houver \alert{linha de visada} entre os dois rádios.

  \medskip
  \begin{itemize}
  \item Os rádios vão \alert{nas pontas}: com quem abre e com quem cerra a fila
  \item Se o grupo se dividir, cada parte leva um rádio
  \item Todos na \alert{na mesma frequência}, combinada antes de sair e gravado na memória
  \end{itemize}

  \medskip
  \begin{atencao}
    Dois rádios lado a lado são inúteis: quem está junto conversa sem rádio.
  \end{atencao}
\end{frame}

\begin{frame}{Comunicação com o mundo}
  \textbf{Frequência Nacional de Chamada (FNC)}: \alert<1>{\freq{146,520}, em simplex}. Chame, estabeleça contato e mude de frequência; fique sempre ouvindo, para casos de emergência.

  \pause
  \medskip
  \textbf{Repetidora}: fica no alto de um morro, recebe na \alert{entrada} e retransmite na
  \alert{saída}. Assim, quem não se enxerga consegue conversar.

  \smallskip
  \centering
  \figura{repetidora}
\end{frame}

\cerjfullimage{caledonia-William}{%
  \usebeamerfont{frametitle}Antenas no Pico do Caledônia\par
  \small Foto de William de Abreu}

\begin{frame}{Exemplos de repetidoras para o PE Três Picos}
  \centering
  \renewcommand{\arraystretch}{1.2}
  \begin{tabular}{@{}lll@{}}
    \toprule
                           & \textbf{PY1MAC }   & \textbf{PY1ABA} \\[1ex]
    Local                  & Pico da Caledônia* & Médio Caledônia* \\
    Saída (você ouve)      & \freq{145,350}     & \freq{147,030} \\
    Entrada (você fala)    & \freq{144,750}     & \freq{147,630} \\
    Deslocamento (offset)  & \freq{-0,6}        & \freq{+0,6} \\
    Subtom                 & nenhum             & \qty{88,5}{\hertz} \\
    \bottomrule
  \end{tabular}

  \medskip
  \raggedright
  \begin{itemize}
  \item Você \alert{ouve na saída} e \alert{fala na entrada}: o rádio troca sozinho, pelo offset
  \item O \alert{subtom} vai junto com a sua voz; sem ele, a PY1ABA não abre
  \end{itemize}
  \vfill\raggedleft
  \comentario{Fonte: RepeaterBook; localização de PY1MAC inferida; PY1ABA no Médio Caledônia segundo listas da LABRE-RJ.}
\end{frame}

% Estudo de alcance: gerado por tools/caledonia.py (relevo Copernicus, picos do OpenStreetMap)
\begin{frame}{Até onde chega a PY1MAC?}
  \centering
  \includegraphics{caledonia-alcance}
\end{frame}

\begin{frame}{Como encontrar repetidoras}
  \begin{columns}
    \column{0.35\textwidth}
    Repetidoras podem ser encontradas no site ou no aplicativo Repeater Book: \linkexterno{\linkRepeaterBook}{repeaterbook.com}.

    \column{0.65\textwidth}
    \foto[height=5cm]{Repeater-Book-RJ}{O Repeater Book permite a descoberta, o cadastro e a atualização de repetidoras. O site também exporta listas prontas para o CHIRP. Na imagem, vê-se a \indicativo{PY1LAG}, repetidora em Teresópolis, RJ.}
  \end{columns}
\end{frame}

\begin{frame}{Antes de sair de casa}
  \begin{atencao}
    \alert{Conheça o seu equipamento} antes de precisar dele! \pause Você sabe operá-lo? Quando tempo dura a sua bateria? Você estã levando uma bateria extra?
  \end{atencao}

  \begin{uncoverenv}<3>
    \begin{atencao}
      Faça a \alert{gestão de riscos}: avise familiares e amigos, combine frequências principais e de \emph{backup} e horários de contato.
    \end{atencao}
  \end{uncoverenv}
\end{frame}

% Protocolo de Áreas Remotas, com base em Salvatti (2024, seção 8.3)
\begin{frame}{Protocolo de operação em áreas remotas}
  Para poupar bateria numa emergência, ligue o rádio só quando há mais chance de alguém ouvir
  \citep{salvatti2024}:

  \medskip
  \begin{enumerate}
  \item Em primeiro lugar, \alert{tente pedir ajuda pelo celular}
  \item Se puder, suba ao ponto mais alto, numa clareira ou com o horizonte desobstruído
  \item Sintonize uma repetidora (melhor cenário), a FNC, \freq{146,520} ou o canal 9 da RUG, \freq{462,6625}
  \item Ligue o rádio nas \alert{horas cheias}, escute por pelo menos 15 minutos e desligue
  \end{enumerate}

  \smallskip
  \comentario{Os 15 minutos são o pior caso: com bateria extra ou power bank, dá para escutar por dias.}
\end{frame}

\begin{frame}{Protocolo de áreas remotas: pedido de socorro}
  \begin{enumerate}
  \setcounter{enumi}{4}
  \item Para interromper uma conversa, diga ``\alert{BREAK}'' três vezes
  \item Ao pedir socorro, mencione \alert{onde, quantas pessoas, o que aconteceu, recursos e suprimentos}. Inicie com:
  \end{enumerate}

  \begin{alertblock}{Pedido de socorro}
    \centering
    ``\textbf{BREAK, BREAK, BREAK! Preciso de ajuda!} Aqui é \textit{\textless identificação\textgreater}, minha localização é \textit{\textless geolocalização\textgreater}.''
  \end{alertblock}

  \begin{enumerate}
  \setcounter{enumi}{5}
  \item Aguarde a resposta
  \end{enumerate}

  \pause
  \begin{atencao}
    Jamais diga ``break'' fora de uma emergência!
  \end{atencao}
\end{frame}

% Radioetiqueta e radioética, com base em Salvatti (2024, p. 17)
\begin{frame}{Radioetiqueta: como se comportar na frequência}
  \begin{enumerate}
  \item Ao sintonizar, \alert<1>{escute primeiro}, fale depois
  \item \alert<2>{Identifique-se} pelo indicativo ao entrar, ao sair e a cada 5 minutos
  \item Seja objetivo; num comunicado longo, faça \alert<3>{pausas de 5 segundos} para outros
    pedirem a palavra
  \item Para entrar numa conversa, \alert<4>{identifique-se na pausa}, peça a palavra e aguarde a sua vez;
    nunca transmita por cima de outra estação
  \item Quem transmite com dificuldade tem \alert<5>{preferência}
  \end{enumerate}

  \begin{uncoverenv}<6->
    \begin{dica}
      Para indicativos e siglas, use o \alert{alfabeto fonético da ITU}:
      \indicativo{PU1JAZ} é \emph{Papa Uniform Um Juliett Alfa Zulu}.
    \end{dica}
  \end{uncoverenv}

  \raggedleft
  \comentario{Com base em \citet[p.~17]{salvatti2024}.}
\end{frame}

\begin{frame}{Radioética: a frequência é pública}
  \begin{itemize}
  \item \alert<1>{Emergência tem a prioridade}: ao ouvir um pedido de socorro, pare de transmitir e
    ajude
  \item \alert<2>{Não use palavrões} ou termos chulos
  \item \alert<3>{Não transmita música}
  \item \alert<4>{Não domine a frequência}: ela é de todos
  \end{itemize}

  \bigskip\raggedleft
  \comentario{Com base em \citet[p.~17]{salvatti2024}.}
\end{frame}

% TODO eventualmente
% \begin{frame}{Diversões}
%   - sota
%   - pota
%   - aprs
%   - Satélites (ISS)
% \end{frame}

%%% Local Variables:
%%% TeX-master: "../presentation"
%%% fill-column: 100
%%% End:
